\subsection{R-17: exact determinant dimension and forced nonrectangular gaps}

\paragraph{Statement R-17.}
The four-variable determinant image has exact affine dimension
\begin{equation}\label{R17:detdimension}
 \dim\mathcal D_{n,4}=2n^2+2\qquad(n\geq1).
\end{equation}
For $n\geq3$ its affine $\mathrm{SL}_4$ quotient consequently
has dimension $2n^2-13$.
Moreover, at each of the padded pairs
\[
 (m,n)=(15,20),\qquad (m,n)=(16,22),
\]
there is a strictly positive multiplicity difference in some
\emph{nonrectangular} partition of exactly four rows in every
sufficiently large coordinate degree. In fact the effective
thresholds $d_{20}$ and $d_{22}$ of G-10 apply. At the second
pair, every positive-degree rectangular multiplicity on the
padded side is zero whenever the rectangular label is integral.

\begin{proof}
At the unperturbed cyclic pencil $M_0$ in R-15, each
complementary matrix has a unique perfect matching.
Therefore its determinant is its permanent multiplied by a
sign. The $(i,j)$ derivatives of the determinant additionally
have the fixed cofactor sign $(-1)^{i+j}$; hence each of the
four corresponding determinant-pencil Jacobian columns is
a sign times the permanent-pencil column. Their full images
are equal to $W_0$. The base-rank proof in R-15 gives
$\dim W_0\geq2n^2+2$ for $n\geq4$.
The determinant image dimension is at least this differential
rank and is at most $2n^2+2$ by G-06. This proves
\eqref{R17:detdimension} in that range. The remaining cases
are the full spaces of linear forms, quadrics, and cubic
surfaces, of dimensions $4,10,20$, respectively, as established
in R-10--R-11 and the elementary linear case.

The stable Fermat point and the $15$-dimensional quotient-fiber
argument of G-11 now give the asserted exact quotient dimension.
This is a direct consequence of the explicit cyclic pencil and
the established group-action bound; no claim of literature
priority for the dimension formula is made.

For the two padded pairs, R-13 and R-14 give
\[
 \begin{array}{c|c|c|c}
 (m,n)&\dim\mathcal P_{m,4}&\dim\mathcal Z_{m,n}
                    &\dim\mathcal D_{n,4}\\ \hline
 (15,20)&\binom{18}{3}=816&819&802\\
 (16,22)&\binom{19}{3}=969&972&970.
 \end{array}
\]
The strict dimension inequalities give exactly-four-row
multiplicity gaps by R-14, with the effective thresholds of
G-10. But $20=4\cdot15/3$, so the equality case of G-12
excludes a positive rectangular gap for the first pair.
For the second, $22>4\cdot16/3$, so the strict case of G-12
makes every positive-degree padded rectangular multiplicity
zero. If $4\nmid nd$, there is no integral four-row rectangle
of size $nd$ in the first place. Thus in every degree covered
by the Hilbert-function argument the required partition is
nonrectangular at both pairs.
\end{proof}

